% Cote 119: offprints, reports, courses, CV and letters (typed and printed part): archive cote 119, site page 166 (PDF page 167).
% typescript. Draft compared: none.
% AI-transcribed by Claude Opus 5.5 (claude-opus-5-5, adaptive thinking, effort high), Message Batches API, batch r119_full (msgbatch_01DwZNm5Y7YsKBWyAN4sFs8w), 2026-11-13; not yet checked by eye.
\scanpage{119/166}
\mspage{- 5 -}

La démonstration est laissée au lecteur qui montrera que l'ensemble
des extrémités des chemins différentiables par morceaux d'origine $a$ est à
la fois ouvert. et fermé dans $U$.

3.2. Soit avec les notations de 3.1. une fonction $F$ continûment différentiable
dans $U$ (les dérivées partielles $\frac{\partial F}{\partial x}$ et $\frac{\partial F}{\partial y}$ existent et sont continues).
Considérons la forme différentielle $\omega = dF$. On a alors évidemment :
(3.2.1.) $\int_{\gamma} dF =$\unsure{;|}$F(b) - F(a)$. Il en résulte que si $dF = 0$ dans $U$
alors $F$ est constante dans $U$.

\uline{Définition 3.2}.

Etant donnée une forme différentielle $\omega$ dans un ouvert connexe $U$
on appelle primitive de $\omega$ une fonction $F$ continûment différentiable dans $U$
telle que $dF = \omega$.

\uline{Remarque} : D'après ce qui précède deux primitives diffèrent d'une constante.
Si $\omega = P\,dx + Q\,dy$ la relation $dF = \omega$ \unsure{équivaut} à
\[ (3.2.2.) \qquad \frac{\partial F}{\partial x} = P \quad \frac{\partial F}{\partial y} = Q \ . \]

\uline{Proposition 3.3.}

Pour que la forme différentielle $\omega$ admette une primitive dans $U$
ouvert connexe de $\mathbb{R}^2$ il faut et il suffit que l'on ait : $\int_{\gamma} \omega = 0$ pour
tout chemin fermé $\gamma$ différentiable par morceaux contenu dans $U$ (dire qu'un
chemin est fermé signifie que son origine et son extrémité sont confondus).

\uline{Démonstration}; La condition est nécessaire d'après (3.2.1.). Supposons que
$\int_{\gamma} \omega = 0$ pour tout chemin fermé dans $U$. Soit $(x_o,y_o) \in U$ ; tout point
$(x,y) \in U$ peut être joint à $(x_o,y_o)$ par un chemin $\gamma$ différentiable par
morceaux. (lemme 3.1.). $\int_{\gamma} \omega$ ne dépend pas de $\gamma$ à cause de notre hypothèse.
Posons $F(x,y) = \int_{\gamma} \omega$ .

\margin{166}
